% !TeX root = ./thuthesis-example.tex

% 论文基本信息配置

% 技术方案基本信息配置（简化版）

% 修复 TimesNewRoman 不支持小型大写字母 (small caps) 的警告
% fontspec 加载 Times New Roman 后内部族名为 TimesNewRoman(0)
% 声明 sc 字形回退到 n，消除 NFSS undefined shape 警告
\AtEndPreamble{%
  \DeclareFontShape{TU}{TimesNewRoman(0)}{m}{sc}{<->ssub*TimesNewRoman(0)/m/n}{}%
  \DeclareFontShape{TU}{TimesNewRoman(0)}{bx}{sc}{<->ssub*TimesNewRoman(0)/bx/n}{}%
}

% 载入所需的宏包

% 定理类环境宏包
\usepackage{amsthm}
% 也可以使用 ntheorem
% \usepackage[amsmath,thmmarks,hyperref]{ntheorem}

% 数学字体配置已简化，使用默认设置

% 可以使用 nomencl 生成符号和缩略语说明
% \usepackage{nomencl}
% \makenomenclature

% 表格加脚注
\usepackage{threeparttable}

% 表格中支持跨行
\usepackage{multirow}
\usepackage{ragged2e}

% 固定宽度的表格。
% \usepackage{tabularx}

% 跨页表格
\usepackage{longtable}
\usepackage{tabularx}
\usepackage{booktabs}  % 提供了 \toprule, \midrule, \bottomrule
\newcolumntype{L}[1]{>{\RaggedRight\arraybackslash}p{#1}}
% 算法
\usepackage{algorithm}
\usepackage{algpseudocode}
\usepackage{listings}
\lstset{
  breaklines=true,
  breakatwhitespace=false,
  columns=fullflexible,
  keepspaces=true
}
%代码字体与颜色
\usepackage{xcolor}
\definecolor{deepblue}{rgb}{0,0,0.5}
\definecolor{deepred}{rgb}{0.6,0,0}
\definecolor{deepgreen}{rgb}{0,0.5,0}
\definecolor{codebg}{RGB}{248,248,248}
\definecolor{termblack}{RGB}{40,44,52}
\definecolor{termgray}{RGB}{220,220,220}

% tcolorbox 代码块和终端环境
\usepackage[most,skins]{tcolorbox}
\IfFontExistsTF{Fira Mono}{%
  \setmonofont{Fira Mono}[Scale=0.82]%
}{%
  \IfFontExistsTF{DejaVu Sans Mono}{%
    \setmonofont{DejaVu Sans Mono}[Scale=0.82]%
  }{}%
}

% 量和单位
\usepackage{siunitx}

% 参考文献使用 BibTeX + natbib 宏包
% 顺序编码制
\usepackage[numbers,sort]{natbib}
\bibliographystyle{thuthesis-numeric}
\usepackage{titlesec}

% 让 \paragraph 参与编号（段级别=4）
\setcounter{secnumdepth}{4}

% 定义 \paragraph 的编号样式为纯阿拉伯数字：1, 2, 3, ...
\makeatletter
\renewcommand\theparagraph{\arabic{paragraph}}
% 如果希望每个 section 内从 1 重新开始编号，保留下一行；
% 如果希望全文连续编号，请删掉下一行。
\@addtoreset{paragraph}{section}
\makeatother

% 把 \paragraph 改成块级标题（标题后换行）
\titleformat{\paragraph}[block]
  {\normalfont\normalsize\bfseries}% 样式
  {\theparagraph}%                  % 显示的编号（比如 1, 2, 3）
  {0.8em}%                          % 编号与标题的间距
  {}                                % 标题前置代码

% 调整段前段后间距（可按需修改）
\titlespacing*{\paragraph}
  {0pt}{3.25ex plus 1ex minus .2ex}{1ex plus .2ex}
\usepackage{indentfirst}  % 首行缩进
\usepackage{fancyhdr}

\AtBeginDocument{%
  \pagestyle{fancy}%
  \fancyhf{}%
  \fancyhead[C]{\songti\zihao{-5} 计算机信息工程学院大作业报告}%
  \fancyfoot[C]{\thepage}%
  \renewcommand{\headrulewidth}{0.4pt}%
  \fancypagestyle{plain}{%
    \fancyhf{}%
    \fancyhead[C]{\songti\zihao{-5} 计算机信息工程学院大作业报告}%
    \fancyfoot[C]{\thepage}%
    \renewcommand{\headrulewidth}{0.4pt}%
  }%
  \fancypagestyle{abstractcn}{%
    \fancyhf{}%
    \fancyhead[C]{\songti\zihao{-5} 计算机信息工程学院大作业报告}%
    \renewcommand{\headrulewidth}{0.4pt}%
  }%
  \fancypagestyle{abstracten}{%
    \fancyhf{}%
    \fancyhead[C]{\rmfamily\zihao{-5} Abstract}%
    \renewcommand{\headrulewidth}{0.4pt}%
  }%
}

% 著者-出版年制
% \usepackage{natbib}
% \bibliographystyle{thuthesis-author-year}

% 生命科学学院要求使用 Cell 参考文献格式（2023 年以前使用 author-date 格式）
% \usepackage{natbib}
% \bibliographystyle{cell}

% 本科生参考文献的著录格式
% \usepackage[sort]{natbib}
% \bibliographystyle{thuthesis-bachelor}

% 参考文献使用 BibLaTeX 宏包
% \usepackage[style=thuthesis-numeric]{biblatex}
% \usepackage[style=thuthesis-author-year]{biblatex}
% \usepackage[style=gb7714-2015]{biblatex}
% \usepackage[style=apa]{biblatex}
% \usepackage[style=mla-new]{biblatex}
% 声明 BibLaTeX 的数据库
% \addbibresource{ref/refs.bib}
%可固定下划线长度
\makeatletter
\newcommand\dlmu[2][4cm]{\hskip1pt\underline{\hb@xt@ #1{\hss#2\hss}}\hskip3pt}
\makeatother

\titleformat{\section}[block]{\Large\bfseries}{\thesection}{1em}{}
% 定义所有的图片文件在 figures 子目录下
\graphicspath{{figures/}}

% 数学命令
\makeatletter
\newcommand\dif{%  % 微分符号
  \mathop{}\!%
  \ifthu@math@style@TeX
    d%
  \else
    \mathrm{d}%
  \fi
}
\makeatother

% 代码块环境：浅色背景 + 行号 + 语法高亮
\newtcblisting{codeblock}[2][]{%
  enhanced, breakable,
  listing only,
  colback=codebg, colframe=codebg!60!black,
  fonttitle=\small\sffamily\bfseries,
  title=#2, attach boxed title to top left={yshift=-2mm, xshift=5mm},
  boxed title style={colback=codebg!60!black, size=small},
  left=4mm, right=3mm, top=4mm, bottom=3mm,
  boxrule=0.6pt, arc=2pt,
  listing options={
    language=#1, style=tcblatex,
    basicstyle=\small\ttfamily,
    keywordstyle=\bfseries\color{deepblue},
    stringstyle=\color{deepgreen},
    commentstyle=\itshape\color{gray},
    numbers=left, numberstyle=\tiny\color{gray},
    numbersep=6pt, stepnumber=1,
    showstringspaces=false,
    breaklines=true, breakatwhitespace=false,
    columns=fullflexible, keepspaces=true,
    tabsize=4, xleftmargin=1em
  }
}

% 终端环境：深色背景 + 仿终端样式
\newtcblisting{terminal}[1][]{%
  enhanced, breakable,
  listing only,
  colback=termblack, colframe=termblack,
  fonttitle=\small\ttfamily,
  title={\color{white}$\triangleright$~#1},
  attach boxed title to top left={yshift=-2mm, xshift=5mm},
  boxed title style={colback=termblack!70!white, size=small},
  left=4mm, right=3mm, top=4mm, bottom=3mm,
  boxrule=0.6pt, arc=2pt,
  listing options={
    language=bash, style=tcblatex,
    basicstyle=\small\ttfamily\color{white},
    keywordstyle=\color{cyan!70!white}\bfseries,
    stringstyle=\color{green!60!white},
    commentstyle=\color{gray}\itshape,
    showstringspaces=false,
    breaklines=true, breakatwhitespace=false,
    columns=fullflexible, keepspaces=true,
    tabsize=4
  }
}

% 保留兼容性：旧 \ttb/\ttm 命令（逐步迁移后可删除）
\newcommand{\ttb}{\bfseries\ttfamily}
\newcommand{\ttm}{\ttfamily}

% 保留兼容性：C++ 代码环境
\newcommand\cppstyle{\lstset{
    language=C++,
    basicstyle=\ttfamily,
    morekeywords={self},
    keywordstyle=\bfseries\color{deepblue},
    emph={MyClass,__init__},
    emphstyle=\bfseries\color{deepred},
    stringstyle=\color{deepgreen},
    frame=tb,
    showstringspaces=false,
    captionpos=t
}}
\lstnewenvironment{cpp}[1][]{\cppstyle\lstset{#1}}{}

% hyperref 宏包在最后调用
\usepackage{hyperref}
